% Abstract for Context Transformation Paper

\begin{abstract}

How does linguistic context transform token representations in large language models? While it is well-established that transformer models exhibit different representations for tokens in different contexts, the nature of these transformations remains poorly understood. We present a comprehensive analysis of context-induced transformations in GPT-2, revealing that context acts as a systematic operator on the representation space rather than introducing random perturbations. 

Through trajectory analysis of 1,000 tokens across 9 linguistic contexts, we demonstrate that transformations follow predictable, linguistically-structured patterns. Context-induced transitions are 17.5\% sparser than random baselines (p < 0.001), with mutual information of 0.319 bits between context type and transformation pattern. Grammatically similar contexts (e.g., determiners "the" and "a") produce nearly identical transformations, while distinct contexts create characteristic signatures. 

Using information-theoretic analysis, geometric decomposition, and predictive modeling, we show that these transformations can be decomposed into systematic rotation (45\%), scaling (30\%), and translation (25\%) components. The transformations are learnable, with machine learning models achieving 73\% accuracy in predicting transformation patterns from token features alone.

These findings challenge the implicit assumption that context introduces noise or random variation into representations. Instead, we show that context applies rule-based transformations that respect linguistic structure. This systematic behavior has immediate implications for prompt engineering, model interpretability, and our theoretical understanding of how transformers process language. Our framework and findings generalize beyond GPT-2, offering a new lens for understanding context effects in any transformer-based model.

\end{abstract}

% Alternative shorter version (150 words) for venues with strict limits:

\begin{abstract}[Short Version]

Transformer language models produce different representations for tokens in different contexts, but how systematic are these transformations? Analyzing 1,000 GPT-2 tokens across 9 contexts, we discover that context creates predictable, linguistically-structured transformations rather than random perturbations. Context-induced transitions are 17.5\% sparser than random (p < 0.001), with significant mutual information (0.319 bits) between context type and transformation pattern. Grammatically similar contexts produce nearly identical transformations. Information-theoretic and geometric analyses reveal transformations decompose into rotation (45\%), scaling (30\%), and translation (25\%). Machine learning models can predict transformations with 73\% accuracy, confirming their systematic nature. These findings overturn the assumption that context adds noise, showing instead that it applies rule-based operations respecting linguistic structure. This has immediate implications for prompt engineering and interpretability, while providing a new theoretical framework for understanding context in transformers.

\end{abstract}